% !PRINTED 123
\weberpage{123}

Ist also $f(b)>g$, so definirt der Schnitt $(\mathfrak A,\mathfrak B)$ eine Zahl $\xi$ im Inneren des Intervalles (1), von der wir nun zeigen können, dass
\begin{equation}
 f(\xi)=g
\tag{3}
\end{equation}
sein muss.

Ist $f(\xi)>g$ und $f(\xi)>g'>g$, so können wir wegen der Stetigkeit von $f(x)$ ein Intervall
\[
 \xi-h\le x\le \xi+h
\]
bestimmen, in dem alle Functionswerthe $f(x)$ grösser als $g'$ und also ihre \gesperrt{untere Grenze} grösser als $g$ ist. Da aber $\xi-h$ zu $\mathfrak A$ gehört, so ist auch in dem ganzen Intervall:
\[
 a\le x\le \xi+h
\]
die \gesperrt{untere Grenze} von $f(x)$ grösser als $g$, während doch $\xi+h$ zu $\mathfrak B$ gehört, worin der Widerspruch liegt.

Es sei jetzt $f(x)$ eine ganze rationale Function mit beliebigen complexen oder reellen Coëfficienten, und die Variable $x$ soll gleichfalls complexe Werthe haben können.

Nach \S.\,31, IV kann, wenn $C$ ein beliebiger positiver Werth ist, die positive Grösse $R$ so angenommen werden, dass
\begin{equation}
 |f(x)|>C
\tag{4}
\end{equation}
ist, sobald
\begin{equation}
 |x|\ge R
\tag{5}
\end{equation}
wird. Zur Veranschaulichung stellen wir das durch die Ungleichung
\[
 |x|<R
\]
begrenzte Gebiet $(R)$ für die Variable $x$ durch eine Kreisfläche vom Radius $R$ in der Ebene der Variablen $x=y+iz$ dar.

Wenn wir $C$ grösser annehmen, als irgend einen Werth von $|f(x)|$ im Inneren des Gebietes $(R)$, zum Beispiel grösser als $|f(0)|$, so wird $|f(x)|$ gewiss im Inneren des Gebietes $(R)$ kleiner werden als an der Begrenzung. Da nun im ganzen Inneren von $(R)$ die Function $|f(x)|$, die wir zur Abkürzung jetzt mit $X$ bezeichnen wollen, nicht negativ wird, so giebt es für die Werthe von $X$ eine \gesperrt{untere Grenze} $g$, und wir haben den Satz zu beweisen:

4. \gesperrt{Es existirt ein Werth $\xi$ von $x$ im Inneren des Gebietes $(R)$, so dass}
\begin{equation}
 |f(\xi)|=g
\tag{6}
\end{equation}
%CONT>>
